\editionPage[\pdfbookmark[0]{Corrigés des défis}{bm-defi-cor-1-parent}]{defi-cor-1}{D01 · Corrigé}{Numéroter tout un livre}{\hyperref[defi-ex-1]{Énoncé p.\,\pageref*{defi-ex-1}}}
\begin{blocEdition}{defi-cor-1}
\editionEtape{Trier par taille}Pages 1 à 9 : 9 pages à un chiffre, soit 9 chiffres. Pages 10 à 99 : 90 pages à deux chiffres, soit 180 chiffres.\editionEtape{Finir le compte}Pages 100 à 150 : 51 pages à trois chiffres, soit 153 chiffres.\editionEtape{Conclure}Total : $9+180+153=\res{342}$ chiffres. Compter des chiffres n’est pas compter des pages : chaque page en demande 1, 2 ou 3.
\end{blocEdition}
\editionPage{defi-cor-2}{D02 · Corrigé}{Le carré magique}{\hyperref[defi-ex-2]{Énoncé p.\,\pageref*{defi-ex-2}}}
\begin{blocEdition}{defi-cor-2}
\editionEtape{La diagonale du 2}$2+5=7$ ; il faut $15-7=8$ dans le coin opposé.\editionEtape{Les autres coins}Avec les pairs aux coins : la première ligne $2+\square+\square$ et les essais organisés donnent la grille 2, 7, 6 ; 9, 5, 1 ; 4, 3, 8.\editionEtape{Contrôler}Chaque ligne, colonne et diagonale totalise $\res{15}$ : par exemple $2+9+4=15$ et $6+1+8=15$.
\end{blocEdition}
\editionPage{defi-cor-3}{D03 · Corrigé}{Plier, plier, plier}{\hyperref[defi-ex-3]{Énoncé p.\,\pageref*{defi-ex-3}}}
\begin{blocEdition}{defi-cor-3}
\editionEtape{Doubler à chaque pli}La suite des épaisseurs est 2, 4, 8, 16, 32 : après 5 plis, $\res{32}$ épaisseurs.\editionEtape{Dépasser 100}On poursuit : 64 après 6 plis, 128 après 7 plis. Il faut $\res7$ plis.\editionEtape{Retenir}Doubler encore et encore fait grandir très vite : bien plus vite qu’ajouter 2 à chaque étape.
\end{blocEdition}
\editionPage{defi-cor-4}{D04 · Corrigé}{Cent œufs, douze boîtes}{\hyperref[defi-ex-4]{Énoncé p.\,\pageref*{defi-ex-4}}}
\begin{blocEdition}{defi-cor-4}
\editionEtape{Un point de départ}Douze boîtes de 6 rangeraient $12\times6=72$ œufs : il en manque 28.\editionEtape{Échanger des boîtes}Chaque échange d’une boîte de 6 contre une boîte de 10 ajoute 4 œufs : $28:4=7$ échanges.\editionEtape{Conclure}Elle utilise $\res{7}$ boîtes de 10 et $\res{5}$ boîtes de 6. Contrôle : $7\times10+5\times6=100$ et $7+5=12$.
\end{blocEdition}
\editionPage{defi-cor-5}{D05 · Corrigé}{La tour triangulaire}{\hyperref[defi-ex-5]{Énoncé p.\,\pageref*{defi-ex-5}}}
\begin{blocEdition}{defi-cor-5}
\editionEtape{Poursuivre la suite}1, 3, 6, 10, 15, 21, 28, 36 : à l’étape 8, $1+2+\cdots+8=\res{36}$ cubes.\editionEtape{Atteindre 45}L’étape 9 ajoute 9 cubes : $36+9=\res{45}$. C’est l’étape 9.\editionEtape{Une astuce}On peut grouper $1+8=2+7=3+6=4+5=9$ : quatre paquets de 9 font 36, sans tout additionner.
\end{blocEdition}
\editionPage{defi-cor-6}{D06 · Corrigé}{Le drapeau étoilé}{\hyperref[defi-ex-6]{Énoncé p.\,\pageref*{defi-ex-6}}}
\begin{blocEdition}{defi-cor-6}
\editionEtape{Rendre les parts égales}En coupant chaque bande en deux, le drapeau compte $3\times2=6$ parts identiques.\editionEtape{Compter}La moitié de la bande du milieu est exactement une de ces 6 parts : les étoiles couvrent $\res{\frac16}$ du drapeau.\editionEtape{Contrôler}Six parts de $\frac16$ reconstituent bien le drapeau entier : $\frac66=1$.
\end{blocEdition}
\editionPage{defi-cor-7}{D07 · Corrigé}{La course des fourmis}{\hyperref[defi-ex-7]{Énoncé p.\,\pageref*{defi-ex-7}}}
\begin{blocEdition}{defi-cor-7}
\editionEtape{Comparer à 1}$\frac78<1$ : C n’a pas atteint l’unité. $\frac54$ et $\frac64$ dépassent 1 : A et B l’ont dépassée.\editionEtape{Départager}Les quarts se comparent entre eux : $\frac64>\frac54$. C’est $\res{\text{B}}$ qui mène, devant A puis C.\editionEtape{Retenir}Comparer à 1 d’abord dégrossit le classement ; les parts égales se comparent ensuite par les numérateurs.
\end{blocEdition}
\editionPage{defi-cor-8}{D08 · Corrigé}{Le nombre au code secret}{\hyperref[defi-ex-8]{Énoncé p.\,\pageref*{defi-ex-8}}}
\begin{blocEdition}{defi-cor-8}
\editionEtape{Mettre en équation}Si le chiffre des unités est $u$, celui des dixièmes est $2\times u$ ; la somme vaut $3\times u=9$.\editionEtape{Conclure}Donc $u=3$ et les dixièmes valent 6 : le nombre est $\res{3{,}6}$.\editionEtape{Contrôler}Le chiffre 6 est bien le double de 3 et $3+6=9$. Un seul nombre convient : la contrainte « trois fois $u$ » n’a qu’une solution.
\end{blocEdition}
\editionPage{defi-cor-9}{D09 · Corrigé}{Entre deux graduations}{\hyperref[defi-ex-9]{Énoncé p.\,\pageref*{defi-ex-9}}}
\begin{blocEdition}{defi-cor-9}
\editionEtape{Passer aux centièmes}Entre $7{,}20$ et $7{,}30$, les écritures à deux décimales sont $7{,}21$, $7{,}22$, …, $7{,}29$ : il y en a $\res9$.\editionEtape{Resserrer}Entre $7{,}25$ et $7{,}30$ : $7{,}26$ à $7{,}29$, soit $\res4$ nombres.\editionEtape{Ouvrir}Avec des rangs supplémentaires, on en trouverait d’autres encore : entre deux nombres différents, il y a toujours de la place.
\end{blocEdition}
\editionPage{defi-cor-10}{D10 · Corrigé}{Cinq euros, huit pièces}{\hyperref[defi-ex-10]{Énoncé p.\,\pageref*{defi-ex-10}}}
\begin{blocEdition}{defi-cor-10}
\editionEtape{Organiser par pièces de 2 euros}Avec deux pièces de 2 euros : il reste 1 euro pour 6 pièces, impossible, car six pièces de 50 c font déjà 3 euros. Avec une pièce de 2 euros : il reste 3 euros pour 7 pièces ; or 7 pièces de 50 c font 3,50 euro, et remplacer une pièce de 50 c par 1 euro augmente encore le total : aucun compte ne tombe sur 3 euros.\editionEtape{Sans pièce de 2 euros}Il faut 5 euros avec 8 pièces de 1 euro ou 50 c. Six pièces de 50 c font 3 euros ; deux pièces de 1 euro complètent : $\res{2\ \text{pièces de 1 euro et 6 pièces de 50 c}}$.\editionEtape{Contrôler}$2\times1+6\times0{,}50=5$ euros et $2+6=8$ pièces : c’est la seule solution.
\end{blocEdition}
\editionPage{defi-cor-11}{D11 · Corrigé}{La corde à border}{\hyperref[defi-ex-11]{Énoncé p.\,\pageref*{defi-ex-11}}}
\begin{blocEdition}{defi-cor-11}
\editionEtape{Le carré}Quatre côtés égaux : $24:4=\res6$ m de côté.\editionEtape{Le rectangle}Deux largeurs font 8 m ; il reste $24-8=16$ m pour deux longueurs : $\res8$ m chacune.\editionEtape{Observer}Même corde, figures différentes : un périmètre ne décide pas de la forme. Contrôle : $2\times(8+4)=24$.
\end{blocEdition}
\editionPage{defi-cor-12}{D12 · Corrigé}{Quatre litres exactement}{\hyperref[defi-ex-12]{Énoncé p.\,\pageref*{defi-ex-12}}}
\begin{blocEdition}{defi-cor-12}
\editionEtape{Premier transvasement}Remplir le 3 L, le verser dans le 5 L : le grand contient 3 L.\editionEtape{Créer 1 L}Remplir à nouveau le 3 L et compléter le grand : il n’accepte que 2 L, il reste $\res{1\ \mathrm L}$ dans le petit.\editionEtape{Finir}Vider le grand, y verser le litre restant, puis un plein bidon de 3 L : $1+3=\res{4\ \mathrm L}$. Chaque étape n’utilise que des bidons pleins ou vides : aucun pifomètre.
\end{blocEdition}
\editionPage{defi-cor-13}{D13 · Corrigé}{Les pièces du carré}{\hyperref[defi-ex-13]{Énoncé p.\,\pageref*{defi-ex-13}}}
\begin{blocEdition}{defi-cor-13}
\editionEtape{Le grand triangle}La moitié de 16 carreaux : $\res8$ carreaux.\editionEtape{Les deux autres pièces}Il reste $16-8=8$ carreaux pour deux pièces égales : $\res4$ carreaux chacune.\editionEtape{Contrôler}$8+4+4=16$ : le découpage recouvre tout le carré, sans trou ni chevauchement. Des formes différentes peuvent avoir la même aire.
\end{blocEdition}
\editionPage{defi-cor-14}{D14 · Corrigé}{L’horloge pressée}{\hyperref[defi-ex-14]{Énoncé p.\,\pageref*{defi-ex-14}}}
\begin{blocEdition}{defi-cor-14}
\editionEtape{De 8 h à 14 h}Six heures s’écoulent : l’avance vaut $6\times3=18$ min. L’horloge affiche $\res{14\ \text h\ 18}$.\editionEtape{Sur un jour entier}En 24 heures, l’avance atteint $24\times3=72$ min, soit $\res{1\ \text h\ 12}$.\editionEtape{Lire à l’envers}À 20 h réelles, 12 heures ont passé : l’avance vaut $12\times3=36$ min et l’horloge affiche justement 20 h 36. Il est donc réellement $\res{20\ \text h}$.
\end{blocEdition}
\editionPage{defi-cor-15}{D15 · Corrigé}{Le trésor du jardin}{\hyperref[defi-ex-15]{Énoncé p.\,\pageref*{defi-ex-15}}}
\begin{blocEdition}{defi-cor-15}
\editionEtape{Construire}Segment de 8 cm, milieu à 4 cm, perpendiculaire à l’équerre, T à 3 cm de M.\editionEtape{Comparer}Le compas pointé en T s’ouvre pareillement vers A et vers B : $\res{TA=TB}$.\editionEtape{Expliquer}En pliant le long de la perpendiculaire en M, A tombe sur B : tout point du pli est à égale distance des deux extrémités. Le trésor est équidistant de A et de B.
\end{blocEdition}
\editionPage{defi-cor-16}{D16 · Corrigé}{Le pliage du carré}{\hyperref[defi-ex-16]{Énoncé p.\,\pageref*{defi-ex-16}}}
\begin{blocEdition}{defi-cor-16}
\editionEtape{La figure obtenue}Le pliage superpose deux moitiés : on obtient un triangle $\res{\text{rectangle et isocèle}}$, avec deux côtés de 8 cm et la diagonale.\editionEtape{Le périmètre}La moitié du périmètre du carré vaut 16 cm, soit deux côtés. Le triangle garde ces deux côtés de 8 cm et ajoute la diagonale : son périmètre $\res{\text{dépasse}}$ 16 cm.\editionEtape{Pourquoi}La diagonale traverse le carré d’un coin à l’autre : elle est plus longue qu’un côté, car le chemin droit entre deux coins opposés dépasse la longueur d’un seul côté.
\end{blocEdition}
\editionPage{defi-cor-17}{D17 · Corrigé}{Les mots du miroir}{\hyperref[defi-ex-17]{Énoncé p.\,\pageref*{defi-ex-17}}}
\begin{blocEdition}{defi-cor-17}
\editionEtape{Deux conditions}Le miroir inverse l’ordre des lettres et retourne chaque lettre. Il faut donc des lettres à axe vertical et un mot identique à l’endroit comme à l’envers.\editionEtape{OTTO}Ses lettres O, T, T, O sont symétriques et le mot renversé reste OTTO : il se lit $\res{\text{inchangé}}$.\editionEtape{Chercher}D’autres exemples : $\res{\text{AVA}}$, TOT, OXO. Un prénom comme ANA échoue : le N n’a pas d’axe vertical.
\end{blocEdition}
\editionPage{defi-cor-18}{D18 · Corrigé}{Le grand cube peint}{\hyperref[defi-ex-18]{Énoncé p.\,\pageref*{defi-ex-18}}}
\begin{blocEdition}{defi-cor-18}
\editionEtape{Les coins}Chaque coin reçoit 3 faces de peinture : il y a $\res8$ coins.\editionEtape{Les arêtes}Entre deux coins, le cube du milieu de chaque arête a 2 faces peintes : 12 arêtes, donc $\res{12}$ cubes.\editionEtape{Le reste}Le centre de chaque face a 1 face peinte (6 cubes) et le cube central n’en a $\res{\text{aucune}}$ : 1 cube. Contrôle : $8+12+6+1=27$.
\end{blocEdition}
\editionPage{defi-cor-19}{D19 · Corrigé}{La photo agrandie}{\hyperref[defi-ex-19]{Énoncé p.\,\pageref*{defi-ex-19}}}
\begin{blocEdition}{defi-cor-19}
\editionEtape{Le lien multiplicatif}De 10 cm à 30 cm, le petit côté a triplé. Le grand côté triple aussi : $15\times3=\res{45}$ cm.\editionEtape{Le cadre}Le tirage mesure 30 cm sur 45 cm, or le cadre n’offre que 40 cm : $45>40$, il $\res{\text{ne convient pas}}$.\editionEtape{Retenir}Un agrandissement multiplie toutes les longueurs par le même nombre : c’est une situation de proportionnalité.
\end{blocEdition}
\editionPage{defi-cor-20}{D20 · Corrigé}{Le tableau effacé}{\hyperref[defi-ex-20]{Énoncé p.\,\pageref*{defi-ex-20}}}
\begin{blocEdition}{defi-cor-20}
\editionEtape{Les totaux de lignes}CM2 : $55-26=29$ élèves.\editionEtape{Case par case}Théâtre CM2 : $29-14=15$. Théâtre CM1 : $24-15=9$. Jeux CM1 : $26-9=17$.\editionEtape{Contrôler}Colonne jeux : $17+14=31$ et $24+31=55$ : tout le tableau est cohérent. Résultat : CM1 : $\res{9}$ théâtre et $\res{17}$ jeux ; CM2 : 15 théâtre et 14 jeux.
\end{blocEdition}
\editionPage{defi-cor-21}{D21 · Corrigé}{Le sac mystère}{\hyperref[defi-ex-21]{Énoncé p.\,\pageref*{defi-ex-21}}}
\begin{blocEdition}{defi-cor-21}
\editionEtape{Traduire « probable »}Plus de la moitié de 10 : au moins $\res6$ boules rouges.\editionEtape{Traduire « possible »}Au moins une bleue : au plus $\res9$ rouges.\editionEtape{Conclure}Le sac contient 6, 7, 8 ou 9 boules rouges : quatre compositions possibles. Les mots du hasard se traduisent en comptes de boules.
\end{blocEdition}
